\section{A lower bound on the tail (Case I), and the scaling iteration (Case II)}
\label{sec:scaling}

This section covers \paper{} (25)--(27). Throughout, $w$ is a cyclic extremal image with
$\tau=\tau(w)\in(0,\frac14]$ and $\varepsilon=\norm{x_0-w}<1$; we write
$m_j=\ip{T^jw}{x_0-w}$, so $m_0=\tau$, $\abs{m_j}\le\min(\tau,\kappa^j)$ and $\norm w\le1$
(Lemma~\ref{lem:defect}), and
\[
\mu=\Bigl(\sum_{j\ge1}\abs{m_j}^2\Bigr)^{1/2}\le\tau\,(L_\tau+2)^{1/2},\qquad
L_\tau=\log_{1/\kappa}(1/\tau).
\]
For a series $c$ with $c_0=1-d$ and $q=dw-\sum_{j\ge1}c_jT^jw$ we have the basic expansion
\begin{equation}\label{eq:basic}
\norm{x_0-V_wc}^2=\varepsilon^2+\norm q^2+2\re\Bigl(d\tau-\sum_{j\ge1}c_jm_j\Bigr).
\end{equation}

\begin{proposition}[Case I]\label{prop:caseI}\stproved{}
Suppose $\abs{m_k}\ge\theta$ for some $k\ge1$, where $0<\theta\le\tau$. Then the extremal series
$b$ of $w$ at level $\varepsilon$ satisfies
\[
\norm b^2\le1-\frac{\theta^2}{8\tau},\qquad
\sum_{j\ge1}\abs{b_j}^2\ge\frac{\theta^4}{256\,\mu^2}\ge\frac{\theta^4}{256\,\tau^2(L_\tau+2)} .
\]
\end{proposition}
\begin{proof}
Let $\eta=\theta/2$, $s=\theta^2/(8\tau)$ and $c=(1-s)e_0+\eta\omega e_k$ with $\abs\omega=1$,
$\omega m_k=\abs{m_k}$. By \eqref{eq:basic},
$\norm{x_0-V_wc}^2\le\varepsilon^2+(s+\eta)^2+2s\tau-2\eta\theta$. Since $s\le\eta/4$,
$(s+\eta)^2\le\frac{25}{16}\eta^2\le\frac32\eta\theta$ and $2s\tau=\frac12\eta\theta$, so $c$ is
feasible. As $\tau\le\frac14$ we have $\eta^2=\theta^2/4\le s/2$ and $s\le\frac1{32}$, so
$\norm c^2=1-2s+s^2+\eta^2\le1-s$.

Now let $b$ be the minimiser, $b_0=1-d$, $\sigma^2=\sum_{j\ge1}\abs{b_j}^2$. From
$\abs{1-d}^2+\sigma^2\le1-s$ we get $\re d\ge s/2$. From \eqref{eq:basic} and feasibility,
$0\ge2\tau\re d-2\sigma\mu$. Hence $\sigma\ge\tau s/(2\mu)=\theta^2/(16\mu)$.
\end{proof}

With $\theta=\tau^4$ this gives $\sum_{j\ge1}\abs{b_j}^2\gtrsim\tau^{14}/L_\tau$, in line with
\paper{} (25). Note that the bound is in terms of the defect of $w$, not of the new extremal
image $V_wb$, whose defect is only bounded above (Proposition~\ref{prop:domination}(iv)).

\subsection*{Case II: the argument of \paper{} is not valid}
Case II is the situation $\abs{m_j}<\tau^4$ for all $j\ge1$. \paper{} takes $\delta=\tau/10$,
lets $w^+$ be the extremal image of $w$ at level $\norm{x_0-(1+\delta)w}$, and claims
\begin{equation}\label{eq:claim27}
\norm{w^+-(1+\delta)w}<\tau^2\quad\text{\paper{} (27)},\qquad\tau(w^+)<(1-\tfrac1{20})\tau ,
\end{equation}
so that iterating gives a norm-convergent sequence with defect tending to $0$.

\begin{remark}[The error]\label{rem:caseIIerror}
The justification of (27) is that lowering $b_0$ below $1+\delta$ must be compensated by
coefficients $b_j$, $j\ge1$, whose effect on the distance is governed by $m_j$, which is tiny.
But the expansion has to be made around $(1+\delta)w$, not around $w$, and
\[
\ip{T^jw}{x_0-(1+\delta)w}=m_j-\delta\ip{T^jw}{w},
\]
whose second term has size $\delta\kappa^j$, far larger than $\tau^4$. Writing
$b_0=1+\delta-s$, $z=\sum_{j\ge1}b_jT^jw$ and $\sigma$ for the tail norm, minimality and
feasibility give
\[
\sigma^2+\abs s^2\le2(1+\delta)\re s,\qquad
\re s\le\frac{\delta^2\norm w^2}{2\tau}+\frac{\sigma\mu}{\tau}
=\frac{\delta\norm w^2}{20}+\frac{\sigma\mu}{\tau},
\]
(the second by minimising $\norm z^2-2\re((s-\delta)\ip wz)$ over all $z\in H$). These allow
$s$ of order $\delta/20$ and a tail $\sigma$ of order $\sqrt\delta$, and give no control of the
sign of $\tau(w^+)-\tau$. Proposition~\ref{prop:27false} shows that $s$ really is of order
$\delta$ for type 1 operators, so that (27) is false.
\end{remark}

\begin{proposition}[(27) is false for type 1 operators]\label{prop:27false}\stproved{}
In Case II with $\tau\le\frac1{10}$, let $j\ge1$, $p=\abs{\ip{T^jw}{w}}$ and $q=\norm{T^jw}$,
and suppose $p\ge45\tau^3$, $q^2\ge8\tau$ and $p^2/q^2\ge100\,\tau\,(L_\tau+2)$. Then
\[
\norm{w^+-(1+\delta)w}\ \ge\ \frac{\tau\,p^2}{100\,q^2}.
\]
In particular, if $T$ is of type 1 and $w/\norm w$ lies in the cone of
Definition~\ref{def:types}, then $p/q\ge2\delta_1\norm w$ for a suitable $j$, and
$\norm{w^+-(1+\delta)w}\ge\frac1{25}\delta_1^2\norm w^2\tau>\tau^2$ as soon as
$\tau<\frac1{25}\delta_1^2\norm w^2$ (and $\tau$ is small enough for the three conditions).
\end{proposition}
\begin{proof}
Let $\abs\omega=1$ with $\omega\ip{T^jw}{w}=p$, $c=\delta p/q^2$, $s=\delta p^2/(40q^2)$ and
$b'=(1+\delta-s)e_0-c\,\omega e_j$. Expanding around $(1+\delta)w$,
\[
\norm{x_0-V_wb'}^2-\norm{x_0-(1+\delta)w}^2
=c^2q^2-2c\delta p+s^2\norm w^2+2scp+2s(\tau-\delta\norm w^2)+2c\re(\omega m_j).
\]
Here $c^2q^2-2c\delta p=-\delta^2p^2/q^2$, $2s\tau=\frac12\delta^2p^2/q^2$,
$2scp\le2s\delta\norm w^2$ (as $p\le q\norm w$), $s^2\norm w^2\le\delta^2p^2/(1600q^2)$ and
$2c\abs{m_j}\le2\delta p\tau^4/q^2$. So the right side is at most
$-0.49\,\delta^2p^2/q^2+2\delta p\tau^4/q^2\le0$, and $b'$ is feasible. Since
$c^2=4\tau s/q^2\le s/2$, $\norm{b'}^2\le(1+\delta)^2-1.4\,s$.

For the minimiser $b$, with $b_0=1+\delta-s^*$, tail norm $\sigma$ and
$\xi=(1+\delta)w-w^+$: minimality gives $2(1+\delta)\re s^*\ge1.4s+\sigma^2$, so
$\re s^*\ge0.6s$ and $\sigma^2\le2.2\re s^*$. Feasibility gives
$0\ge2\tau\re s^*-2\sigma\mu-2\delta\norm\xi\norm w$, that is
$\norm\xi\ge10\re s^*-10\sigma\mu/\tau$. In Case II, $\mu\le\tau^4(L_\tau+2)^{1/2}$, and the
third hypothesis makes $10\sigma\mu/\tau\le\re s^*$. Hence
$\norm\xi\ge9\re s^*\ge5.4\,s=5.4\,\tau p^2/(400q^2)$.
For type 1, some $j\ge1$ has $p\ge\delta_1\norm w^2$, while $q\le\kappa\norm w\le\frac12\norm w$.
\end{proof}

This sharpening of Remark~\ref{rem:caseIIerror} was found by the independent check of these
notes. It concerns only the first claim in \eqref{eq:claim27}; whether the defect decreases is
still open.

\begin{remark}[Historical numerics]
The former script \texttt{case2\_test.py}, preserved under \texttt{experiments/} in Git
snapshot \texttt{7dd4a85}, ran the iteration $w\mapsto w^+$ in the model with $u_1$ a
bump of width $5\cdot10^{-7}$, where $\max_{j\ge1}\abs{m_j}/\tau\approx10^{-4}$. At
$\tau\approx7\cdot10^{-6}$ it finds tail norms $\sigma\approx1.5$--$3\cdot10^{-4}$ (compare
$0.3\sqrt\delta=2.5\cdot10^{-4}$), $\norm{w^+-(1+\delta)w}\approx10^4\,\tau^2$, and ratios
$\tau(w^+)/\tau=0.963,\ 0.985,\ 0.992,\ 0.995,\ 0.996$ over successive steps: the decrease
is not geometric. The strict hypothesis $\abs{m_j}<\tau^4$ cannot be reached in double
precision, so this is not a counterexample to the statement. On re-running, the script also
produced negative computed defects and failed to attain its requested starting residual in
some runs. It has therefore been removed from the maintained experiments; the numbers above
are historical observations, not verified extremal states. The hand argument above does not
depend on them.
\end{remark}

\begin{proposition}[Case II]\label{prop:caseII}\stopen{}
The second claim of \eqref{eq:claim27}, and the consequence that persistent Case II yields a
non-cyclic vector, are unproved. The first claim is false for type 1 operators at small defect
(Proposition~\ref{prop:27false}).
\end{proposition}

\begin{remark}
The exact version is trivial: if $m_j=0$ for all $j\ge1$ then the orbit of $Tw$ is orthogonal to
$x_0-w\ne0$, so $Tw$ is non-cyclic by Lemma~\ref{lem:criterion}. A correct treatment of Case II
should therefore aim at norm convergence with $\sup_{j\ge1}\abs{m_j}\to0$, for which the defect
$\tau=m_0$ need not tend to $0$ at all.
\end{remark}
